\section{Ordered depth generators and exact cancellations}
\label{sec:ordered}

The preceding spectral derivatives retain information about the order of
the summation indices.  A complementary question is whether selected
combinations over \emph{all depths} admit finite evaluations.  We give
such evaluations while keeping one tagged exponent arbitrary.  The
resulting formulas connect ordered harmonic sums to ordinary Hurwitz-zeta
derivatives, generalized Stieltjes constants, and $\log\Gamma$.

The entire interpolation used here is established: its one-endpoint
form is the interpolant of Ihara, Nakamura, and Yamamoto
\cite{INY}, and its relative formulation already appears in the
incoming independent-order report.  Gamma formulas for depth-generating
series have substantial antecedents, including Young's work
\cite{Young}.  The purpose of this section is to derive a common
two-marker formula, prove its convergence and continuation, and extract
explicit cancellations useful for the present research program.

\subsection{Relative ordered harmonic sums}

For $c>0$ and a word $w=(s_1,\ldots,s_d)$, write
\[
 Z_c(w)=\sum_{n_1>\cdots>n_d\ge0}
              \prod_{j=1}^d(n_j+c)^{-s_j},
 \qquad Z_c(\varnothing)=1.
\]
All initial manipulations can be made in the smaller absolute-convergence
domain $\Re s_j>1$.  Let $F_c$ be the formal word series with these
coefficients and constant coefficient one.  Products concatenate words:
the coefficient of $w$ in $FG$ is $\sum_{w=uv}F(u)G(v)$.  Define
\begin{equation}\label{ord:relative}
             \Hrel_{a,b}=F_a^{-1}F_b,\qquad
             h_{a,b}(s)=\zeta(s,b)-\zeta(s,a).
\end{equation}
This specifies every coefficient by a finite algebraic expression.  For
example,
\[
 \Hrel_{a,b}(s)=h_{a,b}(s),\qquad
 \Hrel_{a,b}(s,t)=Z_b(s,t)-Z_a(s,t)-\zeta(s,a)h_{a,b}(t).
\]
The identity $\Hrel_{a,1}(w)=\Psi(w;a-1)$ identifies the one-endpoint
coefficients with \cite[Proposition~2.1]{INY}; their joint entireness in
the spectral indices is \cite[Corollary~4.2]{INY}.  It also proves the
relative assertion without a new regularization prescription, because
\[
                 \Hrel_{a,b}=\Hrel_{a,1}\Hrel_{b,1}^{-1}
\]
and each coefficient of the inverse is a finite polynomial in the
coefficients of $\Hrel_{b,1}$.  We always use these entire spectral continuations.
When $a=b+N$, $N\in\mathbb Z_{\ge0}$, cancellation of the ordered products
gives
\begin{equation}\label{ord:finite}
 \Hrel_{b+N,b}(s_1,\ldots,s_d)
   =\sum_{N>n_1>\cdots>n_d\ge0}
             \prod_{j=1}^d(n_j+b)^{-s_j}.
\end{equation}
Thus this is an interpolation of strict finite harmonic sums, with the
largest index on the left.  Write $1^p$ for a block of $p$ copies of $1$,
including the empty block when $p=0$.

\subsection{A two-marker theorem}

\begin{theorem}[One tagged exponent and two depth markers]\label{ord:two}
Let $a,b>0$, and put $\rho=\min(a,b)$.  The series
\begin{equation}\label{ord:Gdef}
 \mathscr G_s^{a,b}(x,y)
   =\sum_{p,q\ge0}x^py^q\Hrel_{a,b}(1^p,s,1^q)
\end{equation}
converges normally on compact subsets of
$\mathbb C_s\times\{|x|<\rho,\ |y|<\rho\}$.  Define, initially when
$\Re(s+x-y)>1$,
\begin{equation}\label{ord:kernel}
 S_s(c;x,y)=
   \sum_{n\ge0}(n+c)^{1-s}
       \frac{\Gamma(n+c+y)}{\Gamma(n+c+1+x)}.
\end{equation}
Then
\begin{equation}\label{ord:twoformula}
 \boxed{\;
 \mathscr G_s^{a,b}(x,y)
 =\frac{\Gamma(a+x)\Gamma(b)}
        {\Gamma(a)\Gamma(b+y)}
       \bigl(S_s(b;x,y)-S_s(a;x,y)\bigr).
 \;}
\end{equation}
The difference on the right has an entire continuation in $s$, joint
with the marking variables in the stated bidisc.  The completed formula
also has a meromorphic continuation in $x,y$.  On $y=0$ the larger domain
$|x|<a$, with arbitrary $b>0$, is a domain of normal convergence of the
one-marker series.  All spectral derivatives may be taken term by term
on these domains.
\end{theorem}

\begin{proof}
For the finite interval \eqref{ord:finite}, choosing the index $n+b$ of
the tagged exponent and then the subsets above and below it gives
\begin{align}
 \mathscr G_s^{b+N,b}(x,y)
 &=\sum_{n=0}^{N-1}(n+b)^{-s}
       \prod_{k=n+1}^{N-1}\left(1+\frac{x}{k+b}\right)
       \prod_{k=0}^{n-1}\left(1+\frac{y}{k+b}\right) \notag\\
 &=\frac{\Gamma(b+N+x)\Gamma(b)}
          {\Gamma(b+N)\Gamma(b+y)}
    \sum_{n=0}^{N-1}(n+b)^{1-s}
          \frac{\Gamma(n+b+y)}{\Gamma(n+b+1+x)}.
 \label{ord:finitegamma}
\end{align}
Removable Gamma specializations are interpreted through the finite
product in the first line.

Agreement on these finite intervals alone would not justify arbitrary
endpoints.  For that step use triangular matrices
\[
 E_c(n)=
 \begin{pmatrix}
  1+x/(n+c)&(n+c)^{-s}\\
  0&1+y/(n+c)
 \end{pmatrix},
 \qquad P_{c,M}=E_c(M-1)\cdots E_c(0).
\]
Their upper-right entry records exactly one tagged letter, with $x$
counting the letters before it and $y$ those after it.  Consequently
the upper-right entry of $P_{a,M}^{-1}P_{b,M}$ records the corresponding
relative coefficients.

To pass rigorously to exponent $1$, first replace every background
factor $(n+c)^{-1}$ by $(n+c)^{-\tau}$, with real $\tau>1$.
The infinite products then give the coefficients
$\Hrel_{a,b}(\tau^p,s,\tau^q)$ in an absolute-convergence region.
Fix $0<\delta<\rho/2$ and $|x|,|y|<\delta$.  For $\tau\ge1$
sufficiently close to $1$, the logarithms of the diagonal
\emph{quotient} products have differences $O(n^{-\tau-1})$.
After division by the first diagonal entry, each off-diagonal summand
is bounded by
$C(n+1)^{-\Re s+2\delta}$, locally uniformly in the parameters.
Indeed its product factors have logarithmic absolute value at most
$(|x|+|y|)\log(n+1)+O(1)$.
Thus the products and off-diagonal sums converge uniformly for
$\Re s>1+2\delta$ as $\tau\downarrow1$.  Cauchy coefficient extraction
commutes with this limit.  Entireness of the relative coefficients
therefore identifies its coefficients with
$\Hrel_{a,b}(1^p,s,1^q)$.

For completeness, the first and second diagonal entries of $P_{c,M}$
are $A_c(M,x)$ and $A_c(M,y)$, where
\[
 A_c(M,z)=
   \frac{\Gamma(M+c+z)\Gamma(c)}
        {\Gamma(M+c)\Gamma(c+z)}.
\]
Its upper-right entry is
\[
 \frac{\Gamma(M+c+x)\Gamma(c)}
      {\Gamma(M+c)\Gamma(c+y)}
 \sum_{n=0}^{M-1}(n+c)^{1-s}
       \frac{\Gamma(n+c+y)}{\Gamma(n+c+1+x)}.
\]
Inserting these expressions in the inverse of a triangular matrix and
using $\Gamma(M+c+z)/\Gamma(M+c)\sim M^z$ proves
\eqref{ord:twoformula} on the initial region just specified.

It remains to justify the claimed continuation and the convergence
of its Taylor series.  The Gamma-ratio expansion
\cite[\S5.11(iii)]{DLMF} is locally uniform in the markers:
\[
 \frac{\Gamma(t+y)}{\Gamma(t+1+x)}
  =t^{y-x-1}
    \left(\sum_{j=0}^{J-1}c_j(x,y)t^{-j}+O(t^{-J})\right),
 \qquad c_0=1.
\]
The coefficients do not depend on the endpoint; for instance,
$c_1=(y-x-1)(y+x)/2$.  Let $r_J(t)$ be the remainder after multiplying
this expansion by $t^{1-s}$.  Subtraction of the same asymptotic terms
at both endpoints yields
\begin{align}
 S_s(b;x,y)-S_s(a;x,y)
  ={}&\sum_{j=0}^{J-1}c_j(x,y)h_{a,b}(s+x-y+j)\notag\\
    &+\sum_{n\ge0}\bigl(r_J(n+b)-r_J(n+a)\bigr).
 \label{ord:subtract}
\end{align}
The remainder series is normally convergent whenever
$\Re(s+x-y)+J>1$.  Each $h_{a,b}$ is entire: the two simple Hurwitz-zeta
poles have the same residue and cancel.  Increasing $J$ therefore
continues the difference to every $s$, jointly holomorphically on
compact subsets of the stated bidisc.  The finitely many small-$n$
terms present no singularity there.

This continuation has Taylor coefficients in $x,y$ that are entire in
$s$.  They agree with those in \eqref{ord:Gdef} on the initial
half-plane, hence agree everywhere.  The Taylor series of a holomorphic
function on a bidisc converges normally on its compact subsets, proving
the series assertion.  The same subtraction, away from the discrete
Gamma-pole sets, gives meromorphic continuation in the markers.
Finally, at $y=0$ the kernel
$\Gamma(t)/\Gamma(t+1+x)$ is entire in $x$.  The only possible marking
poles of the completed expression then come from $\Gamma(a+x)$.
This proves holomorphy and Taylor convergence for $|x|<a$.
\end{proof}

\subsection{Integer gaps and ordinary zeta derivatives}

The gamma quotient in \eqref{ord:kernel} becomes a polynomial when
the markers have a suitable integer difference.  To state every such
case uniformly, define $c_{m,k}(x)$ by
\[
          (t+x+1)_m=\sum_{k=0}^m c_{m,k}(x)t^k,
 \qquad m\in\mathbb Z_{\ge0},
\]
where $(v)_m=v(v+1)\cdots(v+m-1)$ and $(v)_0=1$.

\begin{corollary}[Finite reduction at an integer marker gap]
\label{ord:gap}
For every complex $s$,
\begin{equation}\label{ord:gapformula}
 \mathscr G_s^{a,b}(x,x+m+1)=
 \frac{\Gamma(a+x)\Gamma(b)}
      {\Gamma(a)\Gamma(b+x+m+1)}
       \sum_{k=0}^m c_{m,k}(x)h_{a,b}(s-1-k).
\end{equation}
This is an ordinary convergent depth sum when the markers belong to the
convergence domain of Theorem~\ref{ord:two}; elsewhere it is an identity
of the specified meromorphic continuations.
In particular, for an integer $m\ge1$, $a>m$, and $b>0$,
\begin{equation}\label{ord:negative}
 \boxed{\;
 \sum_{p\ge0}(-m)^p\Hrel_{a,b}(1^p,s)
 =\frac{\Gamma(a-m)}{\Gamma(a)}
    \sum_{k=0}^{m-1}(-1)^{m-1-k}
           \genfrac{[}{]}{0pt}{}{m}{k+1}h_{a,b}(s-1-k).
 \;}
\end{equation}
Here $\genfrac{[}{]}{0pt}{}{m}{j}$ is an unsigned Stirling number of the
first kind.  The series in \eqref{ord:negative} converges normally
locally in $s$.
\end{corollary}

\begin{proof}
On $y=x+m+1$ the kernel quotient is exactly $(t+x+1)_m$.
Equation \eqref{ord:gapformula} follows first where the defining
Dirichlet series converge, and then everywhere by
\eqref{ord:subtract}.  For \eqref{ord:negative}, set $x=-m$, $y=0$
and use
\[
 \frac{\Gamma(t)}{\Gamma(t+1-m)}
   =\prod_{j=1}^{m-1}(t-j)
   =\sum_{k=0}^{m-1}(-1)^{m-1-k}
                      \genfrac{[}{]}{0pt}{}{m}{k+1}t^k.
\]
The condition $a>m$ places this specialization strictly inside the
one-marker Taylor disc.
\end{proof}

The first two reductions, including their normal-convergence conditions,
are particularly simple:
\begin{align}
 \sum_{p\ge0}(-1)^p\Hrel_{a,b}(1^p,s)
    &=\frac{h_{a,b}(s-1)}{a-1}, && a>1, \label{ord:mone}\\
 \sum_{p\ge0}(-2)^p\Hrel_{a,b}(1^p,s)
    &=\frac{h_{a,b}(s-2)-h_{a,b}(s-1)}
            {(a-2)(a-1)}, && a>2. \label{ord:mtwo}
\end{align}
All their spectral derivatives are finite combinations of ordinary
Hurwitz-zeta derivatives.  At a canceled pole they become Stieltjes
differences.  More explicitly, use the convention
\[
 \zeta(1+u,c)=\frac1u+
       \sum_{k\ge0}\frac{(-1)^k}{k!}\gamma_k(c)u^k.
\]
Differentiating \eqref{ord:mone} at $s=2$ gives, for $a>1$, $b>0$,
\begin{equation}\label{ord:stieltjes}
 \boxed{\;
 \sum_{p\ge0}(-1)^p
    \left.\frac{\partial^k}{\partial s^k}
          \Hrel_{a,b}(1^p,s)\right|_{s=2}
  =\frac{(-1)^k\bigl(\gamma_k(b)-\gamma_k(a)\bigr)}{a-1}.
 \;}
\end{equation}
For $k=0$ this is $(\psi(a)-\psi(b))/(a-1)$.
At the adjacent point $s=1$, the classical relation
$\zeta'(0,c)=\log\Gamma(c)-\tfrac12\log(2\pi)$
\cite[\S25.11]{DLMF} instead yields
\begin{equation}\label{ord:gamma}
 \boxed{\;
 \sum_{p\ge0}(-1)^p
    \left.\frac{\partial}{\partial s}
          \Hrel_{a,b}(1^p,s)\right|_{s=1}
       =\frac{\log\Gamma(b)-\log\Gamma(a)}{a-1}.
 \;}
\end{equation}
For $a=N+1$, $b=1$, $N\ge1$, the sum is finite and its value is
$-\log(N!)/N$.  These statements evaluate ordered combinations over
depth.  They do not remove the nonsymmetric spectral information in
each separate fixed-depth coefficient.

\subsection{A finite-interval formula with several tagged exponents}

\begin{corollary}[A finite staircase of markers]\label{ord:staircase}
Let $a=b+N$, $b>0$, $N\in\mathbb Z_{\ge0}$, and let $r\ge1$ with
$a>r$.  For arbitrary complex $s_1,\ldots,s_r$,
\begin{equation}\label{ord:stairformula}
 \begin{split}
 &\sum_{k_0,\ldots,k_{r-1}\ge0}
       \left(\prod_{j=0}^{r-1}(j-r)^{k_j}\right)
       \Hrel_{a,b}(1^{k_0},s_1,1^{k_1},s_2,\ldots,
                                  1^{k_{r-1}},s_r)\\
 &\hspace{35mm}
       =\frac{\Gamma(a-r)}{\Gamma(a)}
                    \Hrel_{a,b}(s_1-1,\ldots,s_r-1).
 \end{split}
\end{equation}
The displayed sum is finite: all terms of depth greater than $N$ vanish.
\end{corollary}

\begin{proof}
Introduce independent markers $x_0,\ldots,x_r$ for the blocks of ones
before, between, and after the tagged exponents.  The finite-product
calculation in \eqref{ord:finitegamma}, now made at each selected
tagged index $t_j=n_j+b$, gives their joint generator as
\[
 \frac{\Gamma(a+x_0)\Gamma(b)}
      {\Gamma(a)\Gamma(b+x_r)}
 \sum_{N>n_1>\cdots>n_r\ge0}
       \prod_{j=1}^r
           t_j^{1-s_j}
           \frac{\Gamma(t_j+x_j)}
                {\Gamma(t_j+1+x_{j-1})}.
\]
Set $x_j=j-r$.  Every quotient inside the sum becomes one, the last
marker is zero, and the prefactor becomes $\Gamma(a-r)/\Gamma(a)$.
The calculation may first be made with $b>r$, so that no intermediate
Gamma factor is singular.  Both sides are real analytic on the
connected interval $b>\max(0,r-N)$, so the identity extends throughout
that interval.  Equivalently one may use the finite products to fill
all removable intermediate specializations.
\end{proof}

For example, with the same finite-interval hypothesis and $a>2$,
\[
 \sum_{p,q\ge0}(-2)^p(-1)^q
          \Hrel_{a,b}(1^p,s,1^q,t)
       =\frac{\Hrel_{a,b}(s-1,t-1)}{(a-1)(a-2)}.
\]
At $s=t=2$ its value is
\[
 \frac{(\psi(a)-\psi(b))^2-\zeta(2,b)+\zeta(2,a)}
           {2(a-1)(a-2)}.
\]
Corollary~\ref{ord:staircase} is asserted for finite intervals only.
An extension to independent endpoints with several unbounded depth
markers would require its own convergence argument.

\subsection{A polylogarithmic companion}

Use the one-variable multiple polylogarithm convention
\[
 \Li_{s_1,\ldots,s_d}(z)
    =\sum_{n_1>\cdots>n_d\ge1}
        \frac{z^{n_1}}{n_1^{s_1}\cdots n_d^{s_d}},
 \qquad |z|<1.
\]
This convention places the argument on the largest index, as in the
ordered harmonic sums above.

\begin{proposition}[A depth generator and finite-depth evaluations]
\label{ord:polylog}
For $|z|<1$ and arbitrary $s,u\in\mathbb C$, the series
\[
                  \mathscr F_s(z;u)
                    =\sum_{p\ge0}u^p\Li_{1^p,s}(z)
\]
converges normally on compact subsets of these variables and satisfies
\begin{equation}\label{ord:polygen}
 \boxed{\;
 \mathscr F_s(z;u)
   =(1-z)^{-u}
           \int_0^z\frac{(1-t)^u\Li_{s-1}(t)}{t}\,\dd t.
 \;}
\end{equation}
The branch of $\log(1-z)$ vanishes at zero.  In particular,
\begin{align}
 \sum_{p\ge0}(-1)^p\Li_{1^p,s}(z)
   &=(1-z)\bigl(\Li_s(z)+\Li_{1,s-1}(z)\bigr),
       \label{ord:polyone}\\
 \sum_{p\ge0}(-2)^p\Li_{1^p,s}(z)
   &=(1-z)^2\bigl(\Li_s(z)+\Li_{1,s-1}(z)
                                  -\Li_{1,s-2}(z)\bigr)\notag\\
   &\hspace{8mm}+z(1-z)\Li_{s-1}(z).
       \label{ord:polytwo}
\end{align}
More generally every negative integer $u=-m$ admits a finite expression
using ordinary polylogarithms and $\Li_{1,s-j}(z)$ with rational
coefficients in $z$.  All spectral derivatives of these identities
are valid on the same disc.
\end{proposition}

\begin{proof}
For $|u|\le R$ and $s$ in a fixed compact set, the absolute sum of the
coefficients attached to $z^n$ is bounded by
$n^C\prod_{k=1}^n(1+R/k)=O(n^{C+R})$ for a suitable $C$.
Indeed one may first choose the tagged index and then independently
choose the larger background indices.  Summation against $|z|^n$
proves normal convergence, including in $u$ and $s$.

Differentiation of the defining series gives
\[
 \frac{\partial\mathscr F_s}{\partial z}
    =\frac{\Li_{s-1}(z)}z+\frac{u}{1-z}\mathscr F_s(z;u),
 \qquad \mathscr F_s(0;u)=0.
\]
Multiplication by the integrating factor $(1-z)^u$ proves
\eqref{ord:polygen}.  At $u=-m$ use
\[
 \frac1{t(1-t)^m}=\frac1t+\sum_{j=1}^m\frac1{(1-t)^j}.
\]
Thus the integral is $\Li_s(z)+\sum_{j=1}^m I_j(s-1,z)$, where
$I_j(v,z)=\int_0^z\Li_v(t)(1-t)^{-j}\,\dd t$.
Here $I_1(v,z)=\Li_{1,v}(z)$, and integration by parts gives the finite
recursion, for $j\ge2$,
\[
 I_j(v,z)=\frac{\Li_v(z)}{(j-1)(1-z)^{j-1}}
       -\frac{\Li_v(z)+\sum_{k=1}^{j-1}I_k(v-1,z)}{j-1}.
\]
This proves the general finite-depth assertion and, at $m=1,2$,
the two displayed formulas.  Normal convergence permits all derivatives
in the tagged exponent.
\end{proof}

These polylogarithmic formulas, like \eqref{ord:negative}, reduce a
specified depth-generating combination.  They give exact primitives
and spectral derivatives with controlled depth, while leaving open
the separate reduction problems for individual ordered coefficients.
